% ============================================================
\chapter{Calcul Stochastique en Finance}
\label{ch:calcul_stochastique}
% ============================================================

Le mouvement brownien est la brique fondamentale de la finance quantitative --- mais c'est un objet math\'ematiquement monstrueux~: ses trajectoires sont continues mais \emph{nulle part diff\'erentiables}. On ne peut pas d\'eriver au sens classique, et l'int\'egrale ordinaire ne fonctionne plus. C'est Kiyosi It\^o qui, dans les ann\'ees 1940 au Japon, a construit le calcul diff\'erentiel adapt\'e \`a cette situation~: l'\emph{int\'egrale d'It\^o} et la c\'el\`ebre \emph{formule d'It\^o}, qui est au calcul stochastique ce que la r\`egle de la cha\^ine est au calcul classique --- avec un terme correctif suppl\'ementaire, le terme en $\frac{1}{2}\sigma^2 f''$, qui change tout. Ce chapitre construit ces outils et montre comment ils s'appliquent \`a la mod\'elisation financi\`ere.

\section{Mouvement brownien}

\begin{definition}[Mouvement brownien standard]
Un \textbf{mouvement brownien standard} (ou processus de Wiener)
$(W_t)_{t \geq 0}$ est un processus stochastique continu tel que~:
\begin{enumerate}
    \item $W_0 = 0$ p.s.
    \item Les trajectoires $t \mapsto W_t(\omega)$ sont continues p.s.
    \item Les accroissements sont ind\'ependants~: pour $0 \leq s < t \leq u < v$,
          $W_t - W_s$ et $W_v - W_u$ sont ind\'ependants.
    \item Les accroissements sont gaussiens~: $W_t - W_s \sim \mathcal{N}(0, t-s)$.
\end{enumerate}
\end{definition}

\begin{proposition}[Propri\'et\'es du mouvement brownien]
\begin{enumerate}
    \item $\E[W_t] = 0$ et $\Var(W_t) = t$ pour tout $t \geq 0$.
    \item $\Cov(W_s, W_t) = \min(s,t)$.
    \item $W_t$ est une martingale par rapport \`a sa filtration naturelle.
    \item $W_t^2 - t$ est une martingale.
    \item Les trajectoires sont p.s.\ nulle part diff\'erentiables.
    \item La variation quadratique satisfait $\langle W \rangle_t = t$.
\end{enumerate}
\end{proposition}

\begin{codeblock}[title=\textbf{Simulation de trajectoires browniennes}]
\begin{minted}{python}
import numpy as np
import matplotlib.pyplot as plt

def brownian_paths(T, N, M):
    """Simule M trajectoires browniennes sur [0, T] avec N pas."""
    dt = T / N
    dW = np.sqrt(dt) * np.random.standard_normal((M, N))
    W = np.zeros((M, N + 1))
    W[:, 1:] = np.cumsum(dW, axis=1)
    t = np.linspace(0, T, N + 1)
    return t, W

np.random.seed(42)
t, W = brownian_paths(1.0, 1000, 5)
for i in range(5):
    plt.plot(t, W[i])
plt.xlabel('t'); plt.ylabel('$W_t$')
plt.title('Trajectoires browniennes')
plt.grid(True)
plt.show()
\end{minted}
\end{codeblock}

\section{Int\'egrale d'It\^o}

\subsection{Construction}

L'int\'egrale stochastique \'etend l'int\'egrale classique aux int\'egrandes
al\'eatoires par rapport au mouvement brownien.

\begin{definition}[Int\'egrale d'It\^o]
Pour un processus adapt\'e $(H_t)_{t \geq 0}$ v\'erifiant
$\E\!\left[\int_0^T H_t^2\,dt\right] < \infty$, l'\textbf{int\'egrale
d'It\^o} est d\'efinie comme~:
\[
    I_T(H) = \int_0^T H_t\,dW_t
           = \lim_{n \to \infty} \sum_{i=0}^{n-1}
             H_{t_i}(W_{t_{i+1}} - W_{t_i}).
\]
La limite est prise en moyenne quadratique ($L^2$).
\end{definition}

\begin{theoreme}[Propri\'et\'es de l'int\'egrale d'It\^o]
\label{thm:ito_integral_props}
Soit $I_t = \int_0^t H_s\,dW_s$. Alors~:
\begin{enumerate}
    \item $I_t$ est une martingale~: $\E[I_t \mid \mathcal{F}_s] = I_s$
          pour $s \leq t$.
    \item \textbf{Isom\'etrie d'It\^o}~:
          $\E[I_T^2] = \E\!\left[\int_0^T H_t^2\,dt\right]$.
    \item $\E[I_T] = 0$.
    \item La variation quadratique est
          $\langle I \rangle_t = \int_0^t H_s^2\,ds$.
\end{enumerate}
\end{theoreme}

\begin{attention}[title=\textbf{It\^o vs Stratonovich}]
L'int\'egrale d'It\^o \'evalue l'int\'egrande au bord \textbf{gauche}
de chaque intervalle~: $H_{t_i}$ et non $H_{t_{i+1}}$ ou
$\frac{1}{2}(H_{t_i} + H_{t_{i+1}})$. Le choix de Stratonovich
(\'evaluation au milieu) donne un r\'esultat diff\'erent et satisfait
les r\`egles de calcul classiques, mais l'int\'egrale d'It\^o est
privil\'egi\'ee en finance car elle pr\'eserve la propri\'et\'e de
martingale (et donc le pricing risque-neutre).
\end{attention}

\subsection{R\`egles de calcul}

Les r\`egles de calcul diff\`erent du calcul classique. On retient~:
\[
    (dW_t)^2 = dt, \quad (dt)^2 = 0, \quad dW_t \cdot dt = 0.
\]

\section{Formule d'It\^o}

\begin{theoreme}[Formule d'It\^o --- cas unidimensionnel]
\label{thm:ito_formula}
Soit $(X_t)$ un processus d'It\^o satisfaisant
$dX_t = \mu_t\,dt + \sigma_t\,dW_t$
et $f \in C^{1,2}(\R^+ \times \R)$. Alors~:
\[
    df(t, X_t) = \frac{\partial f}{\partial t}\,dt
    + \frac{\partial f}{\partial x}\,dX_t
    + \frac{1}{2}\frac{\partial^2 f}{\partial x^2}\sigma_t^2\,dt.
\]
Sous forme d\'evelopp\'ee~:
\[
    df(t, X_t) = \left(\frac{\partial f}{\partial t}
    + \mu_t\frac{\partial f}{\partial x}
    + \frac{1}{2}\sigma_t^2\frac{\partial^2 f}{\partial x^2}\right)dt
    + \sigma_t\frac{\partial f}{\partial x}\,dW_t.
\]
\end{theoreme}

\begin{formules}[title=\textbf{Formule d'It\^o --- forme compacte}]
\[
    df = \left(\partial_t f + \mu\,\partial_x f
    + \tfrac{1}{2}\sigma^2\partial_{xx} f\right)dt
    + \sigma\,\partial_x f\,dW_t.
\]
Le terme suppl\'ementaire $\frac{1}{2}\sigma^2\partial_{xx}f\,dt$
n'a pas d'analogue en calcul d\'eterministe~; il provient de la
variation quadratique non nulle du brownien.
\end{formules}

\begin{exemple}
Appliquons la formule d'It\^o \`a $f(x) = x^2$ avec $X_t = W_t$~:
\[
    d(W_t^2) = 2W_t\,dW_t + \frac{1}{2}\cdot 2\,dt = 2W_t\,dW_t + dt.
\]
En int\'egrant~: $W_T^2 = 2\int_0^T W_t\,dW_t + T$, d'o\`u~:
\[
    \int_0^T W_t\,dW_t = \frac{W_T^2 - T}{2}.
\]
On note la diff\'erence avec le calcul classique qui donnerait $W_T^2/2$.
\end{exemple}

\begin{theoreme}[Formule d'It\^o multidimensionnelle]
Soient $dX_t^i = \mu_t^i\,dt + \sum_j \sigma_t^{ij}\,dW_t^j$ et
$f \in C^{1,2}$. Alors~:
\[
    df = \partial_t f\,dt + \sum_i \partial_{x_i}f\,dX_t^i
    + \frac{1}{2}\sum_{i,j}\partial_{x_i x_j}f\,
    d\langle X^i, X^j\rangle_t.
\]
\end{theoreme}

\section{Th\'eor\`eme de Girsanov}

\begin{theoreme}[Girsanov]
\label{thm:girsanov}
Soit $(W_t)$ un $\PP$-brownien et $(\theta_t)$ un processus adapt\'e
v\'erifiant la condition de Novikov~:
\[
    \E^\PP\!\left[\exp\!\left(\frac{1}{2}\int_0^T \theta_t^2\,dt\right)\right]
    < \infty.
\]
D\'efinissons la densit\'e de Radon--Nikodym~:
\[
    \frac{d\mathbb{Q}}{d\PP}\bigg|_{\mathcal{F}_T}
    = \mathcal{E}_T(\theta) = \exp\!\left(-\int_0^T \theta_t\,dW_t
    - \frac{1}{2}\int_0^T \theta_t^2\,dt\right).
\]
Alors sous $\mathbb{Q}$, le processus
$\tilde{W}_t = W_t + \int_0^t \theta_s\,ds$
est un mouvement brownien standard.
\end{theoreme}

\begin{intuition}[title=\textbf{Application \`a la finance}]
Le th\'eor\`eme de Girsanov permet de passer de la mesure historique $\PP$
(sous laquelle $dS_t/S_t = \mu\,dt + \sigma\,dW_t$) \`a la mesure
risque-neutre $\mathbb{Q}$ (sous laquelle $dS_t/S_t = r\,dt +
\sigma\,d\tilde{W}_t$) en posant $\theta = (\mu - r)/\sigma$.
La quantit\'e $\theta$ est le \textbf{prix du risque de march\'e}
(\textit{market price of risk}).
\end{intuition}

\section{Th\'eor\`eme de repr\'esentation des martingales}

\begin{theoreme}[Repr\'esentation des martingales]
\label{thm:representation}
Toute martingale de carr\'e int\'egrable $(M_t)_{0 \leq t \leq T}$
adapt\'ee \`a la filtration du brownien peut s'\'ecrire~:
\[
    M_t = M_0 + \int_0^t H_s\,dW_s,
\]
pour un unique processus pr\'evisible $(H_t)$ v\'erifiant
$\E\!\left[\int_0^T H_t^2\,dt\right] < \infty$.
\end{theoreme}

\begin{remarque}
Ce th\'eor\`eme garantit que dans le mod\`ele de Black--Scholes
(march\'e complet), tout actif contingent de carr\'e int\'egrable
est r\'eplicable. Le processus $H_t$ fournit la strat\'egie de couverture.
\end{remarque}

\section{\'Equations diff\'erentielles stochastiques}

\begin{definition}[\'EDS g\'en\'erale]
Une \textbf{\'equation diff\'erentielle stochastique} (EDS) est de la forme~:
\[
    dX_t = b(t, X_t)\,dt + \sigma(t, X_t)\,dW_t, \quad X_0 = x_0,
\]
o\`u $b$ est le coefficient de drift et $\sigma$ le coefficient de diffusion.
\end{definition}

\begin{theoreme}[Existence et unicit\'e]
Si $b$ et $\sigma$ sont lipschitziennes en $x$ uniform\'ement en $t$
et \`a croissance lin\'eaire, alors l'EDS admet une unique solution
forte $(X_t)_{0 \leq t \leq T}$.
\end{theoreme}

\begin{theoreme}[Lien EDP--EDS (Feynman--Kac)]
\label{thm:feynman_kac}
Soit $u(t,x)$ solution de~:
\[
    \partial_t u + b(t,x)\partial_x u
    + \frac{1}{2}\sigma^2(t,x)\partial_{xx}u - r\,u = 0,
    \quad u(T,x) = g(x).
\]
Alors~:
\[
    u(t,x) = \E\!\left[e^{-r(T-t)}g(X_T) \mid X_t = x\right],
\]
o\`u $dX_s = b(s,X_s)\,ds + \sigma(s,X_s)\,dW_s$.
\end{theoreme}

\begin{remarque}
Le th\'eor\`eme de Feynman--Kac \'etablit le lien fondamental entre le
pricing par EDP (Black--Scholes) et le pricing par esp\'erance
risque-neutre (Monte Carlo).
\end{remarque}

\section{Exponentielle stochastique}

\begin{definition}[Exponentielle stochastique]
L'\textbf{exponentielle stochastique} (ou de Dol\'eans-Dade) d'un
processus d'It\^o $X_t$ est~:
\[
    \mathcal{E}_t(X) = \exp\!\left(X_t - X_0 - \frac{1}{2}\langle X\rangle_t\right).
\]
Elle satisfait $d\mathcal{E}_t = \mathcal{E}_t\,dX_t$,
$\mathcal{E}_0 = 1$.
\end{definition}

\begin{proposition}
Le MBG $S_t = S_0\exp\!\left[(\mu - \sigma^2/2)t + \sigma W_t\right]$
est l'exponentielle stochastique de $\mu t + \sigma W_t$~:
\[
    S_t = S_0\,\mathcal{E}_t(\mu\,\cdot + \sigma W).
\]
\end{proposition}

\section{Impl\'ementation~: sch\'emas de discr\'etisation}

\begin{codeblock}[title=\textbf{Sch\'emas d'Euler et de Milstein}]
\begin{minted}{python}
import numpy as np

def euler_maruyama(x0, b, sigma, T, N, M):
    """Schema d'Euler-Maruyama pour dX = b(X)dt + sigma(X)dW."""
    dt = T / N
    X = np.full(M, x0, dtype=float)
    for i in range(N):
        dW = np.sqrt(dt) * np.random.standard_normal(M)
        X = X + b(X)*dt + sigma(X)*dW
    return X

def milstein(x0, b, sigma, sigma_prime, T, N, M):
    """Schema de Milstein (ordre 1.0 fort)."""
    dt = T / N
    X = np.full(M, x0, dtype=float)
    for i in range(N):
        dW = np.sqrt(dt) * np.random.standard_normal(M)
        X = (X + b(X)*dt + sigma(X)*dW
             + 0.5*sigma(X)*sigma_prime(X)*(dW**2 - dt))
    return X

# Application au MBG: dS = r*S*dt + sig*S*dW
r_val, sig = 0.05, 0.20
S0, T, N, M = 100, 1.0, 100, 100000

np.random.seed(42)
ST_euler = euler_maruyama(S0, lambda S: r_val*S, lambda S: sig*S,
                          T, N, M)
ST_milstein = milstein(S0, lambda S: r_val*S, lambda S: sig*S,
                       lambda S: sig, T, N, M)
# Solution exacte
Z = np.random.standard_normal(M)
ST_exact = S0 * np.exp((r_val - 0.5*sig**2)*T + sig*np.sqrt(T)*Z)

print(f"Euler:    E[S_T] = {np.mean(ST_euler):.2f}")
print(f"Milstein: E[S_T] = {np.mean(ST_milstein):.2f}")
print(f"Exact:    E[S_T] = {S0*np.exp(r_val*T):.2f}")
\end{minted}
\end{codeblock}

\section{Exercices}

\begin{exercice}[Formule d'It\^o]
\begin{enumerate}
    \item Calculer $d(e^{W_t})$ et $d(e^{\alpha W_t - \alpha^2 t/2})$.
    \item Montrer que $\mathcal{E}_t = e^{\sigma W_t - \sigma^2 t/2}$ est
          une martingale.
    \item Appliquer la formule d'It\^o \`a $f(t,x) = e^{-rt}x$ pour
          retrouver que $e^{-rt}S_t$ est une $\mathbb{Q}$-martingale.
\end{enumerate}
\end{exercice}

\begin{exercice}[Girsanov]
\begin{enumerate}
    \item V\'erifier explicitement le changement de mesure pour le MBG~:
          \'ecrire $S_t$ sous $\PP$ puis sous $\mathbb{Q}$.
    \item Calculer le prix du risque de march\'e $\theta = (\mu-r)/\sigma$
          pour $\mu = 0.10$, $r = 0.03$, $\sigma = 0.25$.
    \item Montrer que la densit\'e de Radon--Nikodym $\mathcal{E}_T(\theta)$
          est une $\PP$-martingale.
\end{enumerate}
\end{exercice}

\begin{exercice}[Convergence des sch\'emas]
\begin{enumerate}
    \item Comparer l'erreur forte $\E[\abs{S_T^N - S_T}]$ des sch\'emas
          d'Euler et de Milstein en fonction de $N$.
    \item V\'erifier que l'ordre de convergence forte est $1/2$ pour
          Euler et $1$ pour Milstein.
    \item Impl\'ementer un sch\'ema d'ordre faible 2 (Platen) et comparer.
\end{enumerate}
\end{exercice}

\begin{exercice}[Feynman--Kac]
\begin{enumerate}
    \item V\'erifier le th\'eor\`eme de Feynman--Kac pour le call
          europ\'een en comparant la solution de l'EDP de Black--Scholes
          avec l'esp\'erance Monte Carlo.
    \item Appliquer Feynman--Kac \`a l'\'equation de la chaleur
          $\partial_t u = \frac{1}{2}\partial_{xx}u$ et interpr\'eter.
\end{enumerate}
\end{exercice}
